% Part: computability
% Chapter: computability-theory
% Section: halting-problem

\documentclass[../../../include/open-logic-section]{subfiles}

\begin{document}

\olfileid{cmp}{thy}{hlt}
\olsection{থামার সমস্যা}

নির্মাণ অনুসারে, সার্বজনীন আংশিক গণনসাধ্য অপেক্ষক~$\fn{Un}(e,x)$ তখনই
সংজ্ঞায়িত, যখন~$e$ দ্বারা সংকেতায়িত অপেক্ষকটির নিবেশ~$x$-এর গণনা কোনো
মান উৎপন্ন করে। এমনটি ঘটে কি না আমরা নির্ণয় করতে পারি কি না, সে প্রশ্ন
স্বাভাবিক। আসলে তা পারি না। টুরিং-যন্ত্র গণনা-মডেলের ক্ষেত্রে এর অর্থ,
কোনো প্রদত্ত টুরিং যন্ত্র কোনো প্রদত্ত নিবেশে থামে কি না তা গণনার সাহায্যে
নির্ণয় করা যায় না। তাই নিচের উপপাদ্যটি “থামার সমস্যার অনির্ণেয়তা” নামে
পরিচিত। নিচে দুটি প্রমাণ দেব। প্রথমটি আগের আলোচনার ধারাটি চালিয়ে যায়,
দ্বিতীয়টি আরও প্রত্যক্ষ।

\begin{thm}
\ollabel{thm:halting-problem}
ধরা যাক,
\[
h(e, x)  =
\begin{cases}
1 & \text{যদি\/ $\fn{Un}(e, x)$ সংজ্ঞায়িত হয়} \\
0 & \text{অন্যথায়।}
\end{cases}
\]
তাহলে $h$ গণনসাধ্য নয়।
\end{thm}

\begin{proof}
ধরা যাক, $h$ গণনসাধ্য। আমরা দেখাব, এর সাহায্যে একটি সার্বজনীন গণনসাধ্য
অপেক্ষক সংজ্ঞায়িত করা যেত। সংজ্ঞায়িত করি,
\[
\fn{Un'}(e,x) =
\begin{cases}
\fn{Un}(e,x) & \text{যদি $h(e,x) = 1$} \\
0 & \text{অন্যথায়।}
\end{cases}
\]
এখন $\fn{Un'}(e, x)$ একটি সর্বত্র-সংজ্ঞায়িত অপেক্ষক, এবং~$h$ গণনসাধ্য
হলে এটিও গণনসাধ্য। যেমন, আদিম পুনরাবৃত্তি ব্যবহার করে $g$-কে এভাবে
সংজ্ঞায়িত করা যেত:
\begin{align*}
g(0, e, x) & \simeq 0 \\
g(y+1, e, x) & \simeq \fn{Un}(e,x);
\end{align*}
তাহলে
\[
\fn{Un'}(e,x) \simeq g(h(e,x),e,x).
\]
যেখানে $\fn{Un}(e,x)$ সংজ্ঞায়িত, সেখানে $\fn{Un'}(e,x)$ তার সঙ্গে
মেলে। তাই যেসব আংশিক গণনসাধ্য অপেক্ষক সর্বত্র-সংজ্ঞায়িত হয়,
$\fn{Un'}$ তাদের জন্য সার্বজনীন। কিন্তু এটি
\olref[nou]{thm:no-univ}-এর বিরোধী।
\end{proof}

\begin{proof}
ধরা যাক, $h(e,x)$ গণনসাধ্য। অপেক্ষক $g$-কে এভাবে সংজ্ঞায়িত করি:
\[
g(x) =
\begin{cases}
  0                & \text{যদি $h(x,x) = 0$} \\
  \fundefined & \text{অন্যথায়।}
\end{cases}
\]
অপেক্ষক $g$ আংশিক গণনসাধ্য। যেমন, একে
$\umin{y}{h(x,x) = 0}$ হিসেবে সংজ্ঞায়িত করা যায়। অতএব কোনো~$e$-এর
জন্য প্রত্যেক~$x$-এ $g(x) \simeq \fn{Un}(e, x)$। $g$ কি~$e$-এ
সংজ্ঞায়িত? সংজ্ঞায়িত হলে~$g$-এর সংজ্ঞা অনুসারে $h(e,e) = 0$
($g$ সংজ্ঞায়িত হলে তার একমাত্র সম্ভব মান~$0$)। $h$-এর সংজ্ঞা অনুসারে
এর অর্থ $\fn{Un}(e, e)$ অসংজ্ঞায়িত। প্রত্যেক~$x$-এর জন্য
$g(x) \simeq \fn{Un}(e, x)$—এই অনুমান থেকে তখন $g(e)$ অসংজ্ঞায়িত হয়,
যা বিরোধ।
% BN-SRC-145 TARGET-CORRECTION-BEGIN
\par\smallskip\noindent\textnormal{\small\textbf{উৎস-সংশোধন (BN-SRC-145):} হিমায়িত ইংরেজি প্রমাণের বন্ধনীতে বলা হয়েছে, $h$ সংজ্ঞায়িত হলে কেবল~$0$ মান নিতে পারে। কিন্তু অনুমিত~$h$ সর্বত্র-সংজ্ঞায়িত এবং $0$ ও~$1$ দুটিই নিতে পারে; প্রদর্শিত সংজ্ঞা অনুসারে এখানে~$g$-ই সংজ্ঞায়িত হলে কেবল~$0$ নেয়। বাংলা পাঠে সেই অপেক্ষক-নামটি সংশোধিত।} % BN-SRC-145 TARGET-CORRECTION-END
অন্যদিকে, $g(e)$ অসংজ্ঞায়িত হলে $h(e,e) \neq 0$, তাই $h(e,e) = 1$।
ফলে $\fn{Un}(e, e)$ সংজ্ঞায়িত। কিন্তু $g(x) \simeq \fn{Un}(e, x)$
হওয়ায় $g(e)$-ও সংজ্ঞায়িত হতো। আবারও বিরোধ।
\end{proof}

\begin{tagblock}{TMs}
\begin{explain}
এই যুক্তিটি টুরিং যন্ত্রের ভাষায়ও বলা যায়। ধরা যাক, এমন একটি টুরিং
যন্ত্র~$H$ আছে, যা কোনো টুরিং যন্ত্র~$E$-এর বর্ণনা ও একটি নিবেশ~$x$
গ্রহণ করে এবং~$E$ নিবেশ~$x$-এ থামে কি না নির্ণয় করে। তাহলে আমরা আরেকটি
টুরিং যন্ত্র~$G$ বানাতে পারতাম, যা একটিমাত্র নিবেশ~$x$ গ্রহণ করে, সূচক~$x$-এর
যন্ত্র $M_x$ নিবেশ~$x$-এ থামে কি না নির্ণয় করতে $H$ চালায় এবং তার
উল্টো কাজ করে। অর্থাৎ, $H$ যদি জানায় যে $M_x$ নিবেশ~$x$-এ থামে, তবে
$G$ একটি অসীম চক্রে ঢুকে পড়ে; আর $H$ যদি জানায় যে $M_x$ নিবেশ~$x$-এ
থামে না, তবে $G$ সঙ্গে সঙ্গে থামে। নিজের সূচককে নিবেশ হিসেবে পেলে $G$
কি থামে? উপরের যুক্তি দেখায়, সেটি তখনই থামে যখন সেটি থামে না—যা বিরোধ।
অতএব এমন টুরিং যন্ত্র~$H$ আছে—আমাদের এই অনুমান মিথ্যা।
\end{explain}
\end{tagblock}

\end{document}
